% !TEX root = ../nc-cubic.tex
\subsection{偶符号五様体の全関係と重み付き六次式の復元}

\begin{proposition}[重み付き六次式と最小関係]\label{nft:prop:sign-presentation}
命題\ref{nft:prop:sign-generators}の零係数ファイバーで
$H=g(v_1^2,v_2^2,v_3^2,w)$ と置き、$v_i,w$ に重み $1,2$ を付ける。
$H$ は原点に孤立臨界点を持ち、次数を半分にして
\[
 R^{(1)}\simeq\bigoplus_{m\ge0}\operatorname{Jac}(H)_{2m}
\]
となる。$D=(d_{ij})$ を $d_{ii}=u_i$, $d_{ij}=d_{ji}=z_{ij}$ の対称行列とする。
$R$ の最小多項式表示は $D$ の六成分と $w$（次数一）、$P$（次数三）
を変数とし、次の斉次関係を持つ。
\[
 \begin{array}{c|l|c}
  \text{次数}&\text{関係式}&\text{最小本数}\\ \hline
 2&D\text{ の独立な二次小行列式},\quad q&7\\
 3&d_{ij}g_i\quad(1\le i,j\le3)&9\\
 4&Pd_{ij}\quad(1\le i\le j\le3)&6\\
 6&P^2+g_1g_2g_3&1
 \end{array}
\]
この環は socle 次数七の Artin Gorenstein 環で、完全交叉ではない。
同じ偶符号族の滑らかな被覆全体では、$R$ が完全交叉であることも
滑らかな可換 Ku 実現と同値である。
\end{proposition}
\begin{proof}
$S=\C[v_1,v_2,v_3,w]$ に $B$ を $u_i=v_i^2$ で埋め込む。
$H$ の偏微分は $2v_i g_i,q$ である。共通零点があれば
$x_i^2=-g_i(U)$ を解くと、$v_i\ne0$ では $x_i=0$ なので $u_ix_i=0$ である。
これで元の五様体方程式の全偏微分が消え、滑らかさから $U=x=0$、
従って $v=w=0$ となる。よって四偏微分は正則列である。

$S$ は $v_I=\prod_{i\in I}v_i$ を基底とする自由 $B$ 加群である。
偏微分で割ると、偶数個の添字を持つ $v_I$ の係数環は
$I=\varnothing$ で $B/(q,u_lg_l:1\le l\le3)$、
$I=\{i,j\}$ で $B/(q,g_i,g_j,u_k g_k)$ である。
実際、$v_l g_l$ を平方自由語に掛けて $v_I$ の係数を取ればよい。
命題\ref{nft:prop:sign-generators}の全成分を比較すると、
$v_iv_j\mapsto d_{ij}$ は偶数重み部分から $R^{(1)}$ への加群同型で、
小行列式の積の規則により環同型である。

商を取る前の偶数部分は $D,w$ で生成され、六小行列式が関係を尽くす
（命題\ref{nft:prop:quadratic-base}の対分けの証明）。
奇数部分は偶数部分上 $v_1,v_2,v_3$ で生成されるから、
偏微分イデアルの偶数部分は $q$ と九つの $v_jv_i g_i$ で生成される。
$P$ を加えて表の残りの関係で割れば、任意の語は
$R^{(1)}$ の元と $P,wP$ の線型結合に還元できる。ここで
$q(0,0,0,w)$ は非零の $w^2$ の倍数である。これらの直和性は命題\ref{nft:prop:sign-generators}で示した。

最小性も全滑らかな係数で成立する。$q$ は $w^2$ 項を持ち六小行列式から独立である。
九つの三次関係に定数係数の依存が二次関係を法としてあれば、
$S$ 内で $\sum_i L_i(v)v_i g_i=Aq$ が成り立つ。ただし
$L_i$ は $v_j$ の一次結合、$A$ は重み二である。
これは正則列 $(q,v_1g_1,v_2g_2,v_3g_3)$ の重み六の syzygy であり、
最小の Koszul syzygy の重みは $4+5=9$ なので零である。
四次より低い関係の四次への倍数は $P$ を含まず、六つの $Pd_{ij}$ は独立である。
最後の関係は全一次変数を零にすると $P^2$ に移り、それ以前の関係は全て消える。
従って関係イデアルの最小生成元は計 $23$ 本で、
八変数での高さ八を超える。完備局所化してもこの最小本数は変わらない。

重み付き Jacobian 完全交叉の socle 重みは
$4+3\cdot5-(2+3)=14$ である。完全対の偶数部分同士への制限は完全なので、
$R^{(1)}$ の最高次数は七で、その乗法対も完全である。
全環では $P$ と $wP$ が次数三と四に加わる。これらと $R^{(1)}$ の積は
次数七では消える。一方 $P\cdot wP=-w g_1g_2g_3$ は非零である。
これを見るため $B/(q,u_i g_i)\to B/(q,u_i)$ という全射を取る。
前者は上の平方自由基底分解により有限次元であり、後者は
$\C[w]/(w^2)$ である。従って両方とも四元の正則列による商である。
全射を持ち上げる Koszul 複体の射は最上段で $g_1g_2g_3$ 倍であり、
$\operatorname{Ext}^4_B(-,B)$ 上の射はこの元による乗法になる
（次数移動を除く）。有限長加群では低次 Ext が消えるからこの射は単射であり、
後者の $w$ が非零のまま移る。これで $R$ の全次数の乗法対が完全となる。
非零係数では系\ref{nft:cor:allsigns-five}の実核が $R$ を
滑らかな可換三次五様体の Jacobian 完全交叉に同定するので、最後の同値を得る。
\end{proof}

\begin{corollary}[無標識環からの重み付き六次式の復元]\label{nft:cor:sign-sextic}
滑らかな零係数被覆の二つのファイバーで、$R$ が次数付き環として同型であることは、
対応する重み付き六次式 $H$ が重みを保つ多項式座標変換で一致することと同値である。
従ってこの六次式の同型類は残余 dg 圏の不変量である。
\end{corollary}
\begin{proof}
$A=R^{(1)}$ とし、$R/A$ は次数付き $A$ 加群として取る。その次数三から四への
乗法の核
\[
 L=\ker\big(A_1\longrightarrow
          \operatorname{Hom}((R/A)_3,(R/A)_4)\big)
\]
は六成分 $d_{ij}$ の張る超平面である。これは $Pd_{ij}=0$ と
$wP\notin A$ による。$\operatorname{Sym}^2 L\to A_2$ の核は六小行列式の空間である。
実際、残りの二次関係 $q$ は $w^2$ 係数が非零である。従って $L$ の
二次関係から偏極付き Veronese 曲面 $(\mathbb P^2,\mathcal O(2))$ が復元される。
その記述 $L=\operatorname{Sym}^2\langle v_1,v_2,v_3\rangle$ の選択は
三変数の線型変換で尽くされる。射影的には $\mathbb P^2$ の自己同型によるためであり、
全体のスカラーも複素数体では平方根で吸収できる。
$L$ の補方向 $w$ の選択肢は $w\mapsto a w+Q_2(v)$ ($a\ne0$) である。

これで $S=\C[v_1,v_2,v_3,w]$ と、重みを二倍に戻した全射
$S_{\rm even}\to A$ を、重み付き座標変換の違いを除き復元した。
この核 $I_{\rm even}$ から $I=(\partial H)$ の奇数部分も求まる。実際、
\[
 I_{\rm odd}=\{r\in S_{\rm odd}:v_jr\in I_{\rm even}\quad(1\le j\le3)\}.
\]
$J=S/I$ の積の完全対で、$v_j$ 全ての零化イデアルは
$J/(v_1,v_2,v_3)$ の次数 $14$ を基準にした双対である。
この商は $\C[w]/(w^2)$ で偶数重みしか持たない。従って奇数重みの
非零類が全 $v_j$ に消されることはなく、上式が従う。

同じ重み付き勾配イデアル $I$ を与える二つの六次式 $H_0,H_1$ が
重み付き座標変換で移ることを示す。鉛筆 $H_t=(1-t)H_0+tH_1$ の
全偏微分は $I$ に属する。それらが $I/S_+ I$ を張る条件は
各次数で行列が最大階数を持つ条件であり、$0,1$ を含むアフィン直線の開集合を与える。
この上で次数付き Nakayama により $(\partial H_t)=I$ である。
重み付き Euler 式から $H_1-H_0\in I_6$ であり、従って局所的に
\[
 H_1-H_0=\sum_{i=1}^3 b_i(v,t)\partial_{v_i}H_t
                   +c(v,w,t)\partial_w H_t
\]
と書ける。$b_i,c$ は重み一、二で、係数は一定階数の線型方程式の
非零小行列式を選ぶことで $t$ に正則に取れる。有限個の除外点を避けて
$0$ から $1$ に至る区分的滑らかな道を取り、このベクトル場の負の向きに
積分する（道のパラメータではその速度も掛ける）。
$v$ に関する方程式は線型で、$w$ に関する式は $w$ の一次式と
$v$ の二次式の和なので、各小区間の流れは可逆な重み付き多項式変換である。
合成すれば $H_t$ は流れに沿って一定となり、所要の変換を得る。

逆に $H$ の重み付き同型は $A=J_{\rm even}$ と
$L=\langle v_i v_j\rangle$ を保つ。$R$ を回復するには
$A\oplus P\cdot(A/(L))$ に積を戻せばよい。
$P^2$ は $\operatorname{Ann}_A(L)_6$ の非零元である。
この空間は $A$ の socle 次数七の完全対により
$(A/(L))_1^*$ と同型で一次元である。従って二つの $P^2$ の
非零スカラーのずれは $P$ の再尺度化で消え、次数付き環全体が同型となる。
この段階で六次式が分類するのは自然変換環である。
\end{proof}

\begin{theorem}[次数一 del Pezzo 三様体との実比較]\label{nft:thm:sign-sextic-kernel}
零係数の全滑らかな被覆について
\[
 Z_H=V(s^2+H(v,w))\subset\mathbb P(1,1,1,2,3)
\]
は滑らかな指数二、次数一の del Pezzo 三様体であり、
$\mathcal Q\simeq \langle\mathcal O_{Z_H},\mathcal O_{Z_H}(1)\rangle^\perp$
という実 dg Morita 同値がある。特にこの族の二つの全残余圏が同値であること、
その全 Serre 自然変換環が次数付き同型であること、および重み付き六次式 $H$ が
重み付き座標同型であることは、すべて同値である。
\end{theorem}
\begin{proof}
命題\ref{nft:prop:sign-presentation}より $s^2+H$ の affine 臨界点は原点だけである。
$g(0,w)$ は非零の $w^3$ の倍数なので、重み二または三の座標点を通らず、
錐の原点を除いて grading の安定化群は自明である。
従って射影スタックは滑らかな射影スキームで、随伴式は
$\omega_{Z_H}=\mathcal O_{Z_H}(-2)$、
$\mathcal O_{Z_H}(1)^3=6/(1\cdot1\cdot1\cdot2\cdot3)=1$ である。
この次数から $\mathcal O(1)$ は更に可除でなく、指数は二である。
次数付き超曲面環の Gorenstein parameter は $8-6=2$ なので
Orlov \cite[Theorems~2.5(i), 3.10, Proposition~2.17]{NFMOrlovSingularities}の
$i=-1$ の window は $\MF^{\mathbb Z}(s^2+H)$ を例外対
$\mathcal O,\mathcal O(1)$ の左直交に同定する。さらに Serre 関手を合成して
右直交に移す。ここで用いる window は次数切断、導来射影と次数付き
特異圏の dg 商による通常の dg 模型で取る。

$U=(u_1,u_2,u_3,w)$ とし、$g=\sum_j U_j a_j(U)$ と二次式で書く。
$K_g=K(U;a)$ を対応する Koszul 因子化とする。
以下の $T^{\pm1}$ は周期元で、コホモロジー次数は二、内部次数は
源で $-3$、標的で $-6$ である。源の内部次数を二倍して比較する。
源では定理\ref{nfc:thm:counterexample}の生成対象
$E=K_Q\otimes K_g\otimes V$, $Q=\sum u_i x_i^2$ を取る。
三つの $x_i$ の Koszul 収縮と行列の符号の相殺により、
内部シフト間の全 Hom の dg 代数は次の小さい dg 代数と擬同型で結ばれる。
それは $B[T^{\pm1}]$ 上の零微分の代数
\[
 D=\langle\eta_1,\eta_2,\eta_3\mid
 \eta_i\eta_j=\eta_j\eta_i,\ \eta_i^2=-T u_i\rangle
\]
と $\operatorname{End}(K_g)$ の次数付き tensor 積である。
$B=\C[U]$、$\eta_i$ の二重次数は次数倍化後 $(1,-2)$ で、
tensor 積にはコホモロジー次数の Koszul 符号を入れる。
ここで $D$ の表示は通常の積についてである。定理の証明の
$\epsilon^I\theta^J$ 基底での収縮は $B[T^{\pm1}]$ 線型なので、
任意のこの $K_g$ の full endomorphism dg 代数を tensor しても有効である。

標的では $u_i=v_i^2$ と底変更し、$K'=K(s;s)\otimes K_g(v^2,w)$ を取る。
$K(s;s)$ の微分は $s\iota+Ts\epsilon$ であり、
$\theta=\iota-T\epsilon$ は二重次数 $(1,-3)$ の閉演算子で $\theta^2=-T$ である。
$1,\theta$ の代数をその自己 Hom に含めると擬同型になる。
実際、$\epsilon^a\theta^b$ 基底で $[d,\epsilon]=s$ なので、
係数 $s$ の一変数 Koszul 複体が収縮する。
従って標的での全 Hom の小さい dg 模型は
$\C[v,w,T^{\pm1}]\langle\theta\rangle/(\theta^2+T)$ に
同じ $\operatorname{End}(K_g)$ を $B[T^{\pm1}]$ 上で tensor したものである。
後者の内部次数はすべて倍化してある。
標的の偶数内部次数だけを取ると、源の模型からの写像
\[
              u_i\longmapsto v_i^2,\qquad \eta_i\longmapsto v_i\theta
\]
は合成と微分を保つ同型である。実際、左辺の平方自由基底 $\eta^I$ の像は、
可逆な $-T$ の冪を除き、$v_I\theta^{|I|\bmod2}$ であり、
これが右辺偶数部分の全自由基底である。

この比較に必要な厚生成を標的でも証明する。
$\theta$ は $T$ が可逆なので同型を与え、$K'$ の内部奇数シフトは
偶数シフトのコホモロジーシフトから得られる。
そのため偶数シフト全体の厚包の中で、$K'$ に各 $v_i$ の乗法の cone、
すなわち曲率零の Koszul 複体 $K(v_i;0)$ を順に tensor できる。
得る Koszul 因子化の第一列は $(s,w,v_i^2,v_i)$ である。
各 $v_i^2$ から $v_i$ 倍の $v_i$ を引く斉次基底変更により、列は
$(s,w,0,0,0,v_1,v_2,v_3)$ となる。対応する第二列の零スロットの係数は
$a_i(v^2,w)$、$v_i$ スロットの係数は $v_i a_i(v^2,w)$ である。
各 $a_i(v^2,w)$ は $(w,v_1,v_2,v_3)$ に属するので、零スロットの係数をすべて消せる。
具体的に第一列を $A$ とし、第二列に斉次の交代行列と $A$ の積を足す。
零スロットと非零スロットの間の成分で $-a_i$ を表示すれば、非零スロットの
係数は変わらない。この操作は外積二演算子に周期元と対応係数を掛けた
冪零演算子の指数による共役である
（$[\iota_j,\epsilon_l\epsilon_m]=\delta_{jl}\epsilon_m-\delta_{jm}\epsilon_l$）。
従って cone の反復は、$(s,w,v_1,v_2,v_3)$ に対する Koszul 因子化の
各シフトの直和である。これは原点剰余体の stabilization を直和因子に含む。
孤立臨界点では剰余体の全内部シフトが生成する。
定理\ref{nsf:thm:all-dimensions}の生成の証明と同じく、各変数のある冪が
勾配イデアルに入るためその乗法は null-homotopic であり、それらの cone を
順に取ると有限長コホモロジーから元の対象を直和因子として回収できる。

$(3)=[2]$ と $(6)=[2]$ により、源の $E(0),E(1),E(2)$ と
標的の $K'(0),K'(2),K'(4)$ はそれぞれ厚生成する。源側は孤立臨界点での生成と、
ねじれ誘導・平均化および $V$ が唯一の既約射影表現であることによる。
上の Hom の同型と二つの dg 擬同型をこれら三対象間に制限し、
Yoneda と標的生成対象への導来 tensor を合成すれば、実際の Morita 双加群と
逆双加群を得る。これで比較を構成した。
系\ref{nft:cor:sign-sextic}で $R$ から復元した $H$ が一致すれば、
重み付き座標同型（方程式の尺度に応じ $s$ も再尺度化する）のグラフを
両端の比較双加群と合成する。これが最後の十分性を与える。
逆向きは自然変換環の内在性による。
\end{proof}

\begin{corollary}[二つの実現先の内在的判定]\label{nft:cor:sign-full-torelli}
滑らかな偶符号被覆の全残余圏と、滑らかな可換三次五様体の全ての
Kuznetsov 成分の合併を考える。このクラス内では二圏の全次数付き環 $R$ が同型で
あることと、二圏が $\C$-線型 dg Morita 同値であることは同値である。
完全交叉の場合の標的は二次関係を積分した三次五様体であり、
残る場合は上の次数一 del Pezzo 三様体である。
\end{corollary}
\begin{proof}
命題\ref{nft:prop:sign-presentation}と系\ref{nft:cor:allsigns-five}により、
偶符号族で $R$ が完全交叉であることは滑らかな可換三次五様体との
実比較を持つことと同値である。二圏ともこの側にあるとき、
$R_1$ の二次関係は標的方程式の全偏微分の空間である。
命題\ref{gcd:prop:integrability-symmetrizer}により、これを積分した
滑らかな三次式の違いは最細直和成分の非零スカラー倍であり、
複素数体では各成分の座標の三乗根倍で吸収できる。
ゆえに二つの標的は射影同型で、そのグラフと既存の比較核で実同値を得る。
二圏とも非完全交叉の場合は定理\ref{nft:thm:sign-sextic-kernel}を適用する。
二つの場合の環は互いに同型でなく、いずれも逆向きは内在性による。
\end{proof}

標的から自然変換の各次数も照合できる。重み付き MF の Serre 関手は
$(-8)[5]=(-2)[3]$ であり、源の $\rho$ に対応するのは $(2)$ である。
BFK の対角核公式で内部次数 $2m$、射のコホモロジー次数零を取る。
$\mu_6$ の恒等 sector は $\operatorname{Jac}(H)_{2m}$ であり、
原始六乗根では移動空間の余次元五の奇偶条件により消える。
位数三では $s$ のみが固定され、補正後の次数 $2m+5-12$ は
$\operatorname{Jac}(s^2)=\C$ の次数零にならない。
位数二では $w$ のみが固定され、$2m+6-12$ が
$\operatorname{Jac}(g(0,w))=\C[w]/(w^2)$ の次数 $0,2$ を取るのは
$m=3,4$ である。従って源で得た二つの追加成分と全次数で一致する。

\begin{corollary}[次数一三様体側の像]\label{nft:cor:sign-target-locus}
滑らかな Veronese double cone $Z$ の式を
$s^2+w^3+wB_4(v)+B_6(v)=0$ と規格化する。$\Ku(Z)$ が滑らかな偶符号五様体の
全残余圏に同値であるための必要十分条件は、三つの $v$ の共通の線型座標で
$B_4,B_6$ がともに $v_1^2,v_2^2,v_3^2$ の多項式になることである。
この像は次数一三様体のモジュライ内で次元 $13$、余次元 $21$ を持つ。
\end{corollary}
\begin{proof}
直前の六乗根の計算は一般の滑らかな $Z$ にも成立する。
$R_1$ は多項式の重み二の部分であり、これらの積は恒等 sector に留まる。
次数三の追加項があるから $R$ は一次生成でなく、可換三次五様体の
Kuznetsov 成分にはなれない。従って偶符号族との比較があれば零係数点に限る。
定理\ref{nft:thm:sign-sextic-kernel}で $\Ku(Z)\simeq\Ku(Z_H)$ を得る。
Lin--Rennemo--Zhang \cite[Theorem~4.3, Proposition~4.11, Theorem~4.12]{SignLRZVeronese}
は、一般の Veronese double cone について偶数 Jacobian 部分からの復元と
Fourier--Mukai 型同値に対する categorical Torelli を与える。
ここでの dg Morita 同値も射影と包含を合成すれば滑らかな射影多様体間の
perfect 核で表されるのでその仮定を満たす。従って偏極三様体が同型である
（偏極は Picard 群の ample な生成元）。切断環の重み一、二、三の生成元の
変更は $v\mapsto Mv$, $w\mapsto aw+Q_2(v)$,
$s\mapsto bs+Q_3(v,w)$ である。両式の $s$ の一次項と $w^2$ 項が零なので
$Q_3=Q_2=0$ となり、主張の同時座標条件が従う。
逆にこの条件なら $u_i=v_i^2$ で $g$ を定める。偶符号被覆の臨界点から
$v_i^2=u_i$ を選ぶと、$u_ix_i=0$, $x_i^2=-g_i$, $q=0$ により
$H$ の臨界点を得る。$H$ が孤立なら $U=0$、次いで $x_i=0$ となるので、
被覆は滑らかである。上の実比較核が適用できる。

規格化した一般の対 $(B_4,B_6)$ は $15+28=43$ 次元、
固定した符号座標での対は $6+10=16$ 次元を動く。
$M\in\operatorname{GL}_3$ と $a\in\mathbb G_m$ は全体の重み付き dilation を除いて
有効に作用し、群次元は九である。滑らかな sextic $H$ を厳密に保つ
重み付き座標変換は有限である。単位での接空間は四偏微分に対する重み六の
syzygy で、正則列の最初の Koszul syzygy の重みは九だからである。
従って一般の対の軌道次元は九、モジュライ次元は $34$ である。
固定した符号作用を保つ座標変換では $M$ が単項行列である
（三座標点が符号群の共通固定点）。その有効次元は三で、標識付きの次元は $13$。
標識を忘れても、一つの $H$ の射影対称群は有限であり、
移送される射影符号部分群の可能性は有限である。式のスカラーは dilation で
吸収できるのでこの有限性には式を厳密に保つ群を使える。同じ部分群に移る
二表示はその共通固定点を合わせる単項変換で違うだけである。ゆえに像の次元も $13$
で、$34-13=21$ を得る。
\end{proof}

例えば $H=w^3+\sum_i v_i^6$ では六次部分の直線を保つ $v$ の射影変換群は $\Gamma_H=(\mu_6)^2\rtimes S_3$ である。
六次部分の Hessian 行列式が $(v_1v_2v_3)^4$ の非零倍なので、全ての射影変換は
三座標直線を置換し、係数を比べればこの群を得る。
この群内の Klein 四元群の共役類は二つである。置換への像が自明なら
対角群の二捩れ全体である。そうでなければ像は一つの互換であり、
それを $(12)$ と取る。対角核はこの互換が固定する唯一の非自明な
射影符号 $(1,1,-1)$ である。互換の位数二の持上げは射影スカラーを除き
$(a,a^{-1},1)(12)$ ($a^6=1$) で、対角共役で $(12)$ に戻せる。
二つの類は、$w^3$ にそれぞれ
\[
 u_1^3+u_2^3+u_3^3,\qquad
 2u_1^3+30u_1^2u_2+30u_1u_2^2+2u_2^3+u_3^3
\]
を加えた二表示で実現する。後者は $(v_1,v_2)\mapsto(v_1+v_2,v_1-v_2)$
による。どちらの被覆も滑らかである。実際、元の被覆の任意の臨界点から
$v_i^2=u_i$ を選ぶと、$u_i\ne0$ では $x_i=0$ と $g_i=0$ により
$H$ の臨界点を得る。$H$ が孤立なら $U=0$、さらに $x_i^2=-g_i(0)=0$ となる。
二代表の射影符号群は共役でないが、定理\ref{nft:thm:sign-sextic-kernel}により両方の全残余圏は同じ $Z_H$ の残余圏に同値である。

\begin{corollary}[非実現ファイバーの一変数分類]\label{nft:cor:hesse-nonrealizable}
$\lambda=0$, $g_\kappa=w^3+\sum_i u_i^3+\kappa u_1u_2u_3$ とする。
被覆は $\kappa^3\ne-27$ でちょうど滑らかである。この開集合の二点での
全残余 dg 圏が $\C$-線型 dg Morita 同値であることは
$\kappa^3={\kappa'}^3$ と同値である。従ってすべてのファイバーが可換実現障害を
共有しながら、同値類は $\mathbb A^1\setminus\{-27\}$ の複素点で区別される。
\end{corollary}
\begin{proof}
臨界点では $w=0$, $u_i x_i=0$, $x_i^2+3u_i^2+\kappa u_j u_k=0$ となる。
一つでも $u_i$ が零なら全て零であり、全 $u_i\ne0$ では
$x_i=0$ となって Hesse 三次の勾配を得る。その積から
$\kappa^3=-27$ が必要で、$\kappa=-3\xi$, $\xi^3=1$ なら
$(u_1,u_2,u_3)=(1,1,\xi^{-1})$ が臨界点を与える。
非実現は系\ref{nft:cor:allsigns-five}による。

圏同値なら内在的な $R$ が同型なので、系\ref{nft:cor:sign-sextic}により
$w^3+h_\kappa(v)$ が重み付き同型となる。ここで
$h_\kappa=\sum v_i^6+\kappa(v_1v_2v_3)^2$ である。
$w\mapsto a w+Q_2(v)$ の $w^2$ 係数から $Q_2=0$ を得るので、
三変数の線型変換だけを考えればよい。
三次方向微分 $\partial_b^3h_\kappa$ が線型形式の非零三乗になる方向
$[b]$ は、ちょうど三座標点である。実際、方向 $b=(b_1,b_2,b_3)$ の
微分の $v_i^3$, $v_i^2v_j$, $v_1v_2v_3$ の係数は順に
\[
 120b_i^3,\qquad 12\kappa b_j b_k^2,\qquad 48\kappa b_1b_2b_3
 \quad(\{i,j,k\}=\{1,2,3\}).
\]
従って三乗なら $120(\sum_i\eta_i b_i v_i)^3$, $\eta_i^3=1$ と書ける。
二成分だけが非零ならその二変数の $v_i^2v_j$ 係数が一致しない。
三成分とも非零なら巡回する三つの $v_i^2v_j$ 係数の積から
$\kappa^3=30^3$、一方 $v_1v_2v_3$ 係数から $\kappa^3=15^3$ となる。
一成分だけなら確かに非零三乗となる。よって線型変換は座標置換と
対角変換の積であり、純六乗の係数を合わせると、混合係数は三乗根倍のみ変わる。

逆に三乗値が等しければ $r_i^6=1$, $\kappa\prod_i r_i^2=\kappa'$ と選べる。
$u_i\mapsto r_i^2u_i$, $x_i\mapsto r_i^{-1}x_i$, $w\mapsto w$ は
元の被覆の方程式を対応させ、偶符号作用と可換する。同変グラフの核は
選んだ中心文字と例外列を保ち、全残余圏の実同値を与える。
\end{proof}

同じ非実現点を一つの連結な二変数族の中に置ける。実際、
$\lambda x_1x_2x_3$ を戻した被覆の滑らかな底は正確に
\[
 T=\{(\kappa,\lambda):
 (\kappa^3+27)((\kappa\lambda^2+4)^3+27\lambda^6)\ne0\}\subset\mathbb A^2
\]
である。$\lambda\ne0$ では定理\ref{nsf:thm:all-dimensions}の勾配計算を使う。
$u_i=0$ における $g_\kappa$ の制限はいずれも滑らかな Fermat 平面三次であり、
$U=0$ では $\partial_{u_i}F=x_i^2$ が全ての $x_i$ を消す。
残る二条件は $g_\kappa$ および
$g_\kappa+4u_1u_2u_3/\lambda^2$ の滑らかさである。どちらの Hesse 形も
特異なら三つの $u_i$ が全て非零であるから、上の二因子を得る。
$\lambda=0$ は直前の系で確認した。
$T$ は既約であり、定理\ref{nfc:thm:counterexample}の相対例外射影を
この滑らかな射影族と同じ中心文字に適用すれば、$T$ 上の smooth proper 圏族となる。
その可換実現領域は $T\cap\{\lambda\ne0\}$ であり、補集合の圏の同値類を
系\ref{nft:cor:hesse-nonrealizable}が決定する。
